\documentclass[10pt]{article}
\usepackage{lingmacros}
\usepackage{tree-dvips}
\usepackage{makecell}
\begin{document}
\pagenumbering{gobble}
\section*{Energia e fabbisogno}
L'\textbf{energia} (ATP) viene usata per compiere lavoro meccanico e chimico sia a riposo che in attività.\\
Il bilancio energetico è la differenza tra l'apporto energetico (assunto con la dieta) e il dispendio energetico, viene valutato su periodi medio-lunghi per la sua elevata variabilità. Il bilancio può essere positivo, negativo o in equilibrio; il bilancio dev'essere positivo durante gravidanze, crescita e allattamento in condizioni fisiologiche, mentre quelle non fisiologiche sono ad esempio l'anoressia o ustioni (positivo) e l'obesità (negativo). Viene espresso in kcal, mentre il SI lo esprime in kj.\\
L'energia alimentare può essere propria degli alimenti, digeribile e metabolizzabile; quest'ultima è la quota di energia sfruttata per la sintesi di ATP.\\
Il \textbf{dispendio energetico} (DE) su periodi lunghi è il \textbf{DE totale} (DET), il quale corrisponde al valore medio che rappresenta la condizione fisiologica dell'individuo. Il DET è dato dalla somma del \textbf{metabolismo basale} (MB), \textbf{dispendio energetico da attività fisica} (DE-AF), termogenesi da alimenti (TE-AL) e componenti termogeniche minori.
\subsection*{Metabolismo Basale (MB)}
Energia necessaria a mantenere in condizioni basali l'integrità anatomica e funzionale del nostro corpo, è influenzato da età (minore vecchiaia), sesso (maggiore maschi), peso corporeo, ormoni, farmaci e situazioni patologiche. Corrisponde al 50/70\% del DET e la sua maggior parte è destinato a funzioni principali come cervello e cuore.\\
In termini assoluti è kcal/giorno ed è direttamente proporzionale a peso e statura, mentre in termini relativi (per kg di peso) è più elevato nei bambini, più basso nelle donne e aumenta durante il ciclo. Viene valutato in neutralità termica, in posizione supina da 30 min di un individuo vigile, a distanza di 10h dall'ultimo pasto non troppo pesante, rilassamento, assenza di attività fisica intensa nelle ultime 24h: da queste si ricavano formule predittive che permettono di calcolare il MB da variabili \underline{sesso, peso, altezza e età}.
\subsection*{Dispendio Energetico da Attività Fisica (DE-AT)}
Rappresenta qualsiasi aumento di DE dovuto a contrazione muscolare non basale, dipende dall'individuo e dalla quantità di attività fisica sia spontanea che intenzionale; è proporzionale alla massa (quindi MB), non ci sono differenze tra sessi.\\
È indicato in termini assoluti con kcal/min, ma più correttamente è multiplo del MB: livello di attività fisica (LAF) indicato in 24h, rapporto di attività fisica (PAR) moltiplicativo ovvero ogni attività ha un valore per cui viene moltiplicata (camminare= 3MB) e equivalente metabolico di attività (MET) che misura il DE per singola attività in modo ancora più preciso.\\ 
Il \textbf{LAF} è un numero puro che rappresenta la percentuale di DET definendo l'utente inattivo, prevalentemente sedentario, moderatamente attivo, attività fisiche impegnative. I LAF variano per fascia di età, mentre nel caso di gravidanza e allattamento ci si basa sulla differenza tra peso attuale e peso in condizione fisiologica. Per i lattanti non vengono usati i LAF.\\ 
La \textbf{TE-AL} è il dispendio energetico dovuto alla digestione, mentre quelle minori è dovuto a stress, ansia, temperature (TE adattativa)... Hanno ruolo secondario, ma sono significative per fumatori o persone esposte a temperature estreme.\\ 
Il DE viene valutato attraverso calorimetria diretta che misura la liberazione di calore da produzione di ATP e indiretta che può misurare anche MB. Aiutano anche frequenza cardiaca, diario motorio e accelerometri. \\
Secondo i LARN il fabbisogno energetico totale si calcola attraverso \underline{MBxLAF}, mentre per i lattanti (6-12 mesi) si usa DET (formula di Butte) + energia depositata. In utenti in crescita si aggiunge 1\% per la crescita. In gravidanza bisogna aggiungere un quantitativo calorico in base al trimestre, anche l'allattamento al seno richiede un quantitativo aggiuntivo. \\
\begin{center}
\begin{tabular}{ |c|c|c|c| } 
 \hline
 \textbf{} & \textbf{\thead{Fabbisogno energetico\\ totale (DET)}} & \textbf{MB} & \textbf{LAF} \\
 \hline
  \makecell{\textbf{Lattanti} \\\textbf{(6-12 mesi)}} & \makecell{Formula di Butte + \\energia depositata\\
 92,8*peso-152} & &\\
 \hline
  \makecell{\textbf{Bambini} \\\textbf{(1-9 anni)}} & MB*LAF + 1\% & \makecell{59,51*peso-30,4\\ 58,31*peso-31,1\\ \textbf{+3 anni}\\ 22,71*peso+504,3 \\20,32*peso+485,9} & \makecell{1,35-1,39-1,43 \\ \textbf{+3 anni} \\ 1,42-1,57-1,69}\\
  \hline
  \makecell{\textbf{Adolescenti} \\\textbf{(10-17 anni)}} & MB*LAF + 1\% & \makecell{17,69*peso+658,2\\ 13,38*peso+692,6} & 1,66-1,73-1,85 \\
  \hline
  \makecell{\textbf{Adulti} \\\textbf{(18-59 anni)}} & MB*LAF & \makecell{15,06*peso+692,2\\ 14,82*peso+486,6 \\ \textbf{+30 anni} \\ 11,47*peso+873,1\\ 8,13*peso+845,6} & 1,45-1,6-1,75-2,1\\
  \hline
  \textbf{Gravidanza} & \makecell{MB*LAF +\\ DET trimestrale\\ 69 kcal/giorno \\ 266 kcal/giorno\\ 496 kcal/giorno} & &\\
  \hline
  \textbf{Allattamento} & \makecell{MB*LAF + \\500 kcal/giorno\\ (6 mesi)} & &\\
  \hline	
  \makecell{\textbf{Senile} \\ \textbf{(+60 anni)}} & MB*LAF & \makecell{11,71*peso+587,7\\ 9,08*peso+658,5} & 1,4-1,5-1,6-1,75\\
 \hline
\end{tabular}
\end{center}
\section*{Livelli di Assunzione di Riferimento\\ di Nutrimenti ed energia della popolazione\\ italiana (LARN)} 
Si occupano dello studio della popolazione per rappresentare al meglio il fabbisogno medio e l'assunzione adeguata dei vari macronutrienti. L'obiettivo principale dei LARN è l'adeguatezza della dieta, intesa come una dieta che mantenga le capacità chimico-biologiche del nostro organismo in equilibrio e che previene carenze e malattie cronico-degenerative.
\subsection*{Fabbisogno medio (AR)}
Corrisponde al 50\% di un gruppo specifico della popolazione, viene usato per proteine, vitamine e minerali. 
\subsection*{Assunzione di riferimento per la popolazione (PRI)}
Corrisponde al fabbisogno di 97,5\% di soggetti sani in un gruppo specifico della popolazione. Viene usato per proteine, vitamine e minerali. 
\subsection*{Assunzione adeguata (AI)}
Molto simile alla PRI come valore, in base alla situazione viene usata una o l'altra. Viene usato per tutti i nutrienti per lattanti, fibra, vitamine e minerali.
\subsection*{Intervallo di riferimento per assunzione macronutrienti (RI)}
Permette di valutare il rischio associato ad un'assunzione troppo alta o troppo bassa del nutriente, in particolare carboidrati o lipidi. Inoltre permette di fornire una dieta adeguata per quanto riguarda l'introduzione di altri micro e macronutrienti.
\subsection*{Obiettivi nutrizionali per la prevenzione (SDT)}
Viene usato per studiare l'interconnessione tra alimentazione e \textbf{malattie cronico-degenerative}, per questo vengono usati per proteine per individui con più di 60 anni, acidi grassi saturi e colesterolo, ma non vitamine. 
\subsection*{Limite massimo tollerabile di assunzione (UL)}
Corrisponde al valore di assunzione di un alimento che non è associato a effetti avversi basandosi sul concetto di NOAEL, al di sopra di questo valore corrisponde a un rischio potenziale. Viene usato per vitamine e minerali. 
I LARN vengono usati per \textbf{assessment dietetico} che permette di determinare il rischio di inadeguatezza individuale o collettivo, ma anche per\textbf{ planning dietetico} ovvero sviluppare un piano di consumi volto a raggiungere obiettivi nutrizionali individuali o collettivi.
\begin{center}
\begin{tabular}{|c|c|c|}
\hline
\multicolumn{3}{|c|}{\makecell{\textbf{Uso dei LARN nella valutazione }\\\textbf{dello stato di nutrizione e per}\\ \textbf{la sorveglianza nutrizionale}}} \\\hline
& \textbf{A livello individuale} & \textbf{In gruppi di popolazione} \\\hline
\textbf{AR} & \makecell{Si usa (con informazioni sulla\\ variabilità dei fabbisogni e variazione \\intraindividuale degli apporti) per \\esaminare la probabilità che\\ l’apporto usuale sia inadeguato} & \makecell{La proporzione del gruppo con\\ assunzione usuale inferiore all’AR\\ è una stima della prevalenza di \\inadeguatezza} \\\hline
\textbf{PRI} & \makecell{Un’assunzione abituale pari o \\superiore alla PRI si associa a una \\bassa probabilità di inadeguatezza} & \makecell{Da non usare per stimare\\ l’inadeguatezza degli apporti} \\\hline
\textbf{AI} & \makecell{Un apporto abituale pari o\\ superiore all’AI si può assumere\\ come adeguata. Nessuna\\ considerazione può essere fatta se\\ l’introduzione è inferiore all’AI} & \makecell{Un’assunzione media pari \\o superiore all’AI implica una bassa\\ prevalenza di inadeguatezza. Nessuna\\ considerazione  può essere fatta se\\ l’introduzione è inferiore all’AI} \\\hline
\textbf{UL} & \makecell{Un’assunzione abituale \\al di sopra dell’UL \\aumenta il rischio \\di effetti avversi} & \makecell{La proporzione del gruppo con\\ assunzione abituale al di sopra \\dell’UL può considerarsi a rischio di\\ effetti avversi da apporti eccessivi} \\\hline
\multicolumn{3}{|c|}{\makecell{\textbf{Uso dei LARN in dietetica}}} \\\hline
& \textbf{A livello individuale} & \textbf{In gruppi di popolazione} \\\hline
\textbf{AR} & \makecell{Non utilizzare come obiettivo\\ di introduzione. Questo\\ livello si associa ad una\\ probabilità di inadeguatezza di 50\%} & \makecell{Ridurre al minimo la proporzione \\di popolazione con apporti al di \\sotto dell’AR. In questo caso \\l’apporto medio risulterà\\ probabilmente superiore alla PRI.} \\\hline
\textbf{PRI} & \makecell{Mirare a questo livello \\di apporto per rendere minima \\la probabilità di inadeguatezza} & \makecell{Considerare come il livello \\minimo d’assunzione del \\nutriente che va garantito} \\\hline
\textbf{AI} & \makecell{Garantire questo livello di \\assunzione per minimizzare la \\probabilità di inadeguatezza} & \makecell{Pianificare un’assunzione media\\ pari all’AI per ridurre al\\ minimo il rischio d’inadeguatezza} \\\hline
\textbf{UL} & \makecell{Mirare ad un apporto abituale\\ al di sotto dell’UL per evitare \\rischi di effetti avversi} & \makecell{Ridurre al minimo la proporzione\\ del gruppo con introduzione al di \\sopra dell’UL per escludere il rischio\\ di effetti avversi} \\\hline
\textbf{STD}& \multicolumn{2}{|c|}{\makecell{Gli STD vanno considerate sia per l'individuo che per la \\ collettività a livello quantitativo e qualitativo per prevenire\\ il rischio di malattie cronico degenerative.}}\\\hline
\end{tabular}
\end{center}
\section*{Macro e micronutrienti}
Sono contenuti negli alimenti e sono indispensabili per il nostro organismo dove compiono funzioni strutturali, energetiche e bioregolatrici.\\
I \textbf{macronutrienti} vengono scissi in unità monomeriche e sono: carboidrati e fibra, lipidi, proteine e acqua.\\
I \textbf{micronutrienti} sono introdotti in piccole quantità ($<1g/100g$ di alimento crudo) e non vengono modificati come vitamine e sali minerali.\\
L'\underline{etanolo} è un non nutriente che fornisce un quantitativo di energia. \\
\subsection*{Carboidrati e fibre}
Costituiti da catene di CHO principalmente vegetali tranne il lattosio.
Sono monosaccaridi, disaccaridi, oligosaccaridi (3-9) e polisaccaridi($>10$).\\
Hanno \underline{funzione energetica} fornendo mediamente 4kcal/g, il glucosio in particolare 3,75kcal/g. Hanno \underline{azione-antichetogena} impedendo la formazione di corpi chetonici con un minimo apporto di 100-130 g/giorno. 
\paragraph*{Disponibilità nutrizionale}Quelli disponibili possono essere usati direttamente dalle cellule, mentre altri come fibra e polialcoli non sono digeribili e vengono solo fermentati dal microbiota intestinale. L'\textbf{amido} basa la sua digeribilità sulla grandezza (inversamente proporzionale) e la presenza di amilopectina (direttamente proporzionale) \\
\\
Gli \underline{zuccheri aggiunti} sono inseriti nella preparazione oltre a quelli naturali.\\
L'\textbf{indice glicemico} (IG) classifica gli alimenti in base all'incremento glicemico dopo un pasto predendo da riferimento glucosio (IG 100) o pane bianco. La risposta glicemica varia anche dalla persona e da fattori alimentari come: tipo e forma dello zucchero, metodi di preparazione dell'alimento e quantità di altri nutrienti e fibre. \\
Il \textbf{carico glicemico} (GL) quantifica l'effetto glicemico in base alla quantità e la qualità del carboidrato, in particolare è usato nei pazienti diabetici. Viene calcolato moltiplicando la quantità di carboidrati in 100g per il suo IG diviso 100.\\
\subsubsection*{Alimenti con carboidrati}
Il \textbf{glucosio} è contenuto in forma libera o polimero in frutta, verdura e dolcificanti naturali (miele).\\
Il \textbf{fruttosio} è contenuto in miele e frutta. Quello in commercio si ottiene dall'amido di mais.\\ 
Il \textbf{galattosio} non si trova in forma libera ma solo sottoforma di lattosio.\\
Il \textbf{saccarosio} raffinato è zucchero bianco, mentre se contiene melassa è detto bruno. Viene estratto dalla barbabietola da zucchero e dalle canne. \\
Il \textbf{lattosio} è l'unico carboidrato animale, è presente nel latte e derivati. Per essere digerito è necessaria la \underline{lattasi} e a livello industriale viene usato per preparazioni a fresco e come conservante. \\
Il \textbf{maltosio} è presente naturalmente nei semi germogliati e viene usato per birra e gelatine di frutta.\\
L' \textbf{amido} è principalmente in pane, pasta, riso, cereali, legumi e radici.
\subsubsection*{Fibra} 
Composta da molecole resistenti a idrolisi e assorbimento raggiungendo il colon senza modifiche. È vegetale e collabora nella fermentazione intestinale.\\
Quella \textbf{solubile} è contenuta nella polpa di frutta e legumi, stimola la crescita del microbiota e riduce l'assorbimento di glucosio, lipidi e colesterolo.\\
Quella \textbf{insolubile} è contenuta nella buccia della verdura, aiuta nella defecazione ammorbidendo le feci e velocizzando il processo pulendo l'intestino. 
\begin{center}
\begin{tabular}{|c|c|c|c|}
\hline
&\makecell{\textbf{STD}\\\textbf{Obiettivo nutrizionale}} & \makecell{\textbf{RI}\\\textbf{Intervallo di}\\\textbf{riferimento per}\\\textbf{l'assunzione di}\\\textbf{macronutrienti}} & \makecell{\textbf{AI}\\\textbf{Assunzione}\\\textbf{adeguata}} \\\hline
\makecell{\textbf{Carboidrati}\\\textbf{totali} \\\textbf{(CHOt)}} & \makecell{Prediligere fonti alimentari\\ amidacee a basso GI in\\ particolare quando gli apporti \\di carboidrati disponibili si\\ avvicinano al limite superiore \\dell'RI. Limitare alimenti\\ in cui la riduzione del GI\\ è ottenuta aumentando il\\ consumo in fruttosio o lipidi}& 45-60\% Energia & \\\hline
\makecell{\textbf{Zuccheri}\\\textbf{(CHOs)}} & \makecell{Limitare il consumo di\\ zuccheri a $<15\%$ energia, un \\apporto totale $>25\%$ energia è\\ potenzialmente legato a eventi \\avversi sulla salute. Limitare\\ alimenti ad alto contenuto di\\ fruttosio perché può \\causare diarrea.} & nd & nd\\\hline
\makecell{\textbf{Fibra}\\\textbf{alimentare}} & \makecell{Preferire alimenti naturalmente\\ ricchi in fibra alimentare quali \\cereali integrali, legumi, frutta\\ e verdura. Negli adulti,\\ consumare almeno 25g/giorno\\ di fibra alimentare, anche in\\ caso di apporti energetici\\ $<2000kcal/giorno$.} & \makecell{Adulti:\\12,6-16,7g/1000kcal} & \makecell{Età evolutiva:\\8,4g/1000kcal}\\\hline
\end{tabular}
\end{center}
\subsection*{Proteine}
\end{document}

